%!TEX root = main.tex
In recent years, there have emerged many works on resource management for different types of networks. Table \ref{tab:related_work} summarizes these works, from aspects of network scenarios (w.r.t. network types, blockages, user mobility), resource management variables, objectives and techniques adopted.

Among all these works, traffic routing and transmission scheduling problems have been carefully studied
\cite{polese2018distributed, kilpi2017link,saad2019millimeter}.
However, only a few consider different link statuses (i.e., line-of-sight (LOS), non-LOS (NLOS), outages, etc.).
The works \cite{garcia2018delay,niu2019relay} find routing paths to bypass links interrupted by obstacles.
The work in \cite{he2015minimum} performs a slot-by-slot link scheduling to maximize the instantaneous throughput considering link blockage behaviors described by discrete-time Markov chains.
Authors in \cite{hu2017relay} deal with the relay selection and link scheduling problem to maximize the end-to-end throughput, using 3D models of buildings as primary blockage sources.
In \cite{yao2017outage}, heuristic algorithms for user scheduling and power allocation are proposed to reduce outage occurrences.
Our work differs from the above works in two aspects. First, we focus on randomly mobile 3D obstacles that can pose new challenges to the resource management problem, compared with static obstacles or stochastic-process-modeled blockages. Second, we aim to train intelligent transmission schemes that can make proper decisions by implicitly predicting the future moments.

User mobility in mmWave networks is mostly addressed in handover / user-cell association \cite{khan2019reinforcement,khosravi2020learning}, and only a few other types of resource management problems consider it.
The work in \cite{kim2019online} proposes a contextual multi-armed bandit algorithm to schedule transmissions to users with unknown positions.
In \cite{tang2020deep}, deep Q-network (DQN) is exploited to allocate capacity for up/down link transmissions in a 5G heterogeneous network with high-mobility.
Authors in \cite{pateromichelakis2018context} select routing paths based on the mobility and traffic conditions, then they perform a link-resource time sharing according to flow occupations.
A particle swarm optimization (PSO) algorithm is described in \cite{tafintsev2020aerial} to properly manage unmanned aerial vehicles (UAVs) in mmWave aerial access and backhaul networks to serve mobile users.
In \cite{shen2020mobility}, a band and beam allocation scheme for mmWave networks is presented, which considers massive multiple-input multiple-output (MIMO) systems and user trajectories.

However, none of these works deal with user mobility in coordinating concurrent backhaul and access transmission in mmWave IAB networks that are also characterized by link blockages, node buffers and different duplexing modes.


Recent years have witnessed a widespread utilization of RL techniques in mmWave networks.
A large part of the works in the literature deal with throughput maximization.
The work in \cite{lei2020deep} defines a spectrum allocation for IAB networks that maximizes the sum of log-rates through double DQN and actor critic techniques.
Authors in \cite{vu2018path} resort to regret RL and successive convex approximation to perform route selection and rate allocation, respectively.
Risk-sensitive RL is adopted in \cite{vu2018ultra} to control transmitter beamwidth and power so as to maximize data rate.
In \cite{zhang2021resource}, the authors propose a resource allocation framework based on advantage actor critic and column generation to maximize the throughput of static mmWave IAB networks.
Some of the other works investigate latency performances. Link scheduling approaches based on deep deterministic policy gradient (DDPG) \cite{gupta2020learning} and multi-armed bandit \cite{ortiz2019scaros} are proposed to minimize the end-to-end latency in mmWave backhaul networks.
Interference management and capacity issues have been faced as well. The work in \cite{feng2019dealing} allocates capacity between the core network and mmWave BSs to users subject to blockages. Authors in \cite{elsayed2020transfer} mitigate the inter-beam inter-cell interference through joint user-cell association and selection of number of beams.

The works \cite{vu2018path,zhang2021resource,gupta2020learning,ortiz2019scaros} study the path routing and link scheduling, however, \cite{zhang2021resource} did not consider UE mobility, while \cite{vu2018path,gupta2020learning} did not consider both link blockages and UE mobility.
Though \cite{ortiz2019scaros} considered both, it focuses on backhaul operations, emphasizing the load dynamics implicitly affected by user mobility statistics. Similarly, \cite{gupta2020learning} focuses on the backhaul part of an IAB network, while assuming static nodes and no blockages.


Finally, MARL emerges as a promising approach to cope with traffic signal control \cite{chu2019multi} and 
resource allocation \cite{kwon2019multi} in vehicular networks, user association \cite{sana2020multi} and handover management \cite{guo2020joint} in mmWave networks.
In \cite{kwon2019multi}, the authors propose a multi-agent deep deterministic policy gradient (MADDPG)-based approach for vehicles to select BSs to associate with and channels to communicate on, in order to maximize the system revenue.
The authors in \cite{sana2020multi} provide an MARL framework for static UEs to be associated with BSs, based on the hysteretic deep recurrent Q-network (HDRQN) algorithm.
In \cite{guo2020joint}, the authors jointly handle handover and power allocation to improve throughput and reduce handover frequency, by developing an MARL algorithm based on proximal policy optimization (PPO).
These works demonstrate good performance of MARL algorithms in making decisions in sophisticated systems, however, none of them deal with resource allocation problems considered in this work.


The above works set good examples of performing resource management in mmWave networks. However, there is still a lack of approaches that can maximize system throughput by adaptively performing flow allocation and link scheduling for both access and backhaul parts of mmWave IAB networks, whose dynamics are caused by both intermittent links and UE mobility. This work is motivated by such issues. 
In our previous work \cite{zhang2021mobility}, we have considered a simplified sector-based blockage model, a star backhaul topology, and only HD IAB-nodes. In this article, we extend it by considering a realistic link blockage model based on 3D mobile obstacles, a general tree-like backhaul topology, and both FD and HD working modes for IAB-nodes.


\begin{table*}[]
\centering
\caption{Summary of related work.}
\label{tab:related_work}
\resizebox{\textwidth}{!}{%
\begin{tabular}{|cc|ccccccc|cccc|}
\hline
\multicolumn{2}{|c|}{\multirow{2}{*}{}} &
  \multicolumn{7}{c|}{Resource management} &
  \multicolumn{4}{c|}{Objective} \\ \cline{3-13} 
\multicolumn{2}{|c|}{} &
  \multicolumn{1}{c|}{\begin{tabular}[c]{@{}c@{}}Routing/\\ path selection/\\ relay control\end{tabular}} &
  \multicolumn{1}{c|}{\begin{tabular}[c]{@{}c@{}}Transmission\\ scheduling\end{tabular}} &
  \multicolumn{1}{c|}{\begin{tabular}[c]{@{}c@{}}Capacity/\\ rate\\ allocation\end{tabular}} &
  \multicolumn{1}{c|}{\begin{tabular}[c]{@{}c@{}}Spectrum/\\ channel\\ allocation\end{tabular}} &
  \multicolumn{1}{c|}{\begin{tabular}[c]{@{}c@{}}Beam\\ allocation\end{tabular}} &
  \multicolumn{1}{c|}{\begin{tabular}[c]{@{}c@{}}Power\\ allocation\end{tabular}} &
  \begin{tabular}[c]{@{}c@{}}User-cell\\ association/\\ handover\end{tabular} &
  \multicolumn{1}{c|}{\begin{tabular}[c]{@{}c@{}}Throughput/\\ rate\\ maximization\end{tabular}} &
  \multicolumn{1}{c|}{\begin{tabular}[c]{@{}c@{}}Latency/\\ delay\\ minimization\end{tabular}} &
  \multicolumn{1}{c|}{\begin{tabular}[c]{@{}c@{}}Hops/\\ time length\\ minimization\end{tabular}} &
  \begin{tabular}[c]{@{}c@{}}Network\\ utility\\ maximization\end{tabular} \\ \hline
\multicolumn{2}{|c|}{IAB networks} &
  \multicolumn{1}{c|}{\begin{tabular}[c]{@{}c@{}}\cite{polese2018distributed,vu2018path,zhang2021resource}\\ \textbf{Ours}\end{tabular}} &
  \multicolumn{1}{c|}{\begin{tabular}[c]{@{}c@{}}\cite{zhang2021resource}\\ \textbf{Ours}\end{tabular}} &
  \multicolumn{1}{c|}{\cite{vu2018path}} &
  \multicolumn{1}{c|}{\cite{lei2020deep}} &
  \multicolumn{1}{c|}{} &
  \multicolumn{1}{c|}{} &
   &
  \multicolumn{1}{c|}{\begin{tabular}[c]{@{}c@{}}\cite{polese2018distributed,lei2020deep,zhang2021resource}\\ \textbf{Ours}\end{tabular}} &
  \multicolumn{1}{c|}{} &
  \multicolumn{1}{c|}{\cite{polese2018distributed}} &
  \cite{vu2018path} \\ \hline
\multicolumn{2}{|c|}{\begin{tabular}[c]{@{}c@{}}mmWave \\ backhaul networks\end{tabular}} &
  \multicolumn{1}{c|}{\cite{garcia2018delay,niu2019relay,hu2017relay,ortiz2019scaros}} &
  \multicolumn{1}{c|}{\begin{tabular}[c]{@{}c@{}}\cite{kilpi2017link,saad2019millimeter,garcia2018delay,niu2019relay,hu2017relay}\\ \cite{ortiz2019scaros,gupta2020learning}\end{tabular}} &
  \multicolumn{1}{c|}{} &
  \multicolumn{1}{c|}{} &
  \multicolumn{1}{c|}{} &
  \multicolumn{1}{c|}{} &
   &
  \multicolumn{1}{c|}{\cite{saad2019millimeter,garcia2018delay,niu2019relay, hu2017relay}} &
  \multicolumn{1}{c|}{\cite{kilpi2017link,garcia2018delay,gupta2020learning,ortiz2019scaros}} &
  \multicolumn{1}{c|}{\cite{kilpi2017link}} &
   \\ \hline
\multicolumn{2}{|c|}{Other networks} &
  \multicolumn{1}{c|}{\cite{he2015minimum,pateromichelakis2018context,tafintsev2020aerial}} &
  \multicolumn{1}{c|}{\begin{tabular}[c]{@{}c@{}}\cite{he2015minimum,yao2017outage}\\ \cite{kim2019online,pateromichelakis2018context}\end{tabular}} &
  \multicolumn{1}{c|}{\cite{tang2020deep,feng2019dealing}} &
  \multicolumn{1}{c|}{\cite{kwon2019multi}} &
  \multicolumn{1}{c|}{\begin{tabular}[c]{@{}c@{}}\cite{khosravi2020learning,shen2020mobility}\\ \cite{vu2018ultra,elsayed2020transfer}\end{tabular}} &
  \multicolumn{1}{c|}{\begin{tabular}[c]{@{}c@{}}\cite{yao2017outage,vu2018ultra}\\ \cite{guo2020joint}\end{tabular}} &
  \begin{tabular}[c]{@{}c@{}}\cite{khan2019reinforcement,sana2020multi,elsayed2020transfer}\\ \cite{kwon2019multi,khosravi2020learning,guo2020joint}\end{tabular} &
  \multicolumn{1}{c|}{\begin{tabular}[c]{@{}c@{}}\cite{he2015minimum,khan2019reinforcement,khosravi2020learning}\\ \cite{kim2019online,tang2020deep,pateromichelakis2018context,tafintsev2020aerial,shen2020mobility,feng2019dealing}\\ \cite{elsayed2020transfer,sana2020multi,guo2020joint}\end{tabular}} &
  \multicolumn{1}{c|}{} &
  \multicolumn{1}{c|}{} &
  \cite{vu2018ultra,kwon2019multi} \\ \hline
\multicolumn{2}{|c|}{Blockages} &
  \multicolumn{1}{c|}{\begin{tabular}[c]{@{}c@{}}\cite{garcia2018delay,niu2019relay,he2015minimum,hu2017relay,zhang2021resource,ortiz2019scaros}\\ \textbf{Ours}\end{tabular}} &
  \multicolumn{1}{c|}{\begin{tabular}[c]{@{}c@{}}\cite{garcia2018delay,niu2019relay,he2015minimum,hu2017relay,yao2017outage}\\ \cite{zhang2021resource,ortiz2019scaros}\\ \textbf{Ours}\end{tabular}} &
  \multicolumn{1}{c|}{\cite{feng2019dealing}} &
  \multicolumn{1}{c|}{} &
  \multicolumn{1}{c|}{\cite{shen2020mobility,vu2018ultra}} &
  \multicolumn{1}{c|}{\cite{yao2017outage,vu2018ultra}} &
   &
  \multicolumn{1}{c|}{\begin{tabular}[c]{@{}c@{}}\cite{garcia2018delay,niu2019relay,he2015minimum,hu2017relay,shen2020mobility,zhang2021resource}\\ \cite{feng2019dealing}\\ \textbf{Ours}\end{tabular}} &
  \multicolumn{1}{c|}{\cite{garcia2018delay,ortiz2019scaros}} &
  \multicolumn{1}{c|}{\cite{yao2017outage,he2015minimum}} &
  \cite{vu2018ultra} \\ \hline
\multicolumn{2}{|c|}{UE mobility} &
  \multicolumn{1}{c|}{\begin{tabular}[c]{@{}c@{}}\cite{pateromichelakis2018context,tafintsev2020aerial,ortiz2019scaros}\\ \textbf{Ours}\end{tabular}} &
  \multicolumn{1}{c|}{\begin{tabular}[c]{@{}c@{}}\cite{yao2017outage,kim2019online}\\ \cite{pateromichelakis2018context,ortiz2019scaros}\\ \textbf{Ours}\end{tabular}} &
  \multicolumn{1}{c|}{\cite{tang2020deep}} &
  \multicolumn{1}{c|}{\cite{lei2020deep,kwon2019multi}} &
  \multicolumn{1}{c|}{\begin{tabular}[c]{@{}c@{}}\cite{khosravi2020learning,shen2020mobility}\\ \cite{elsayed2020transfer}\end{tabular}} &
  \multicolumn{1}{c|}{\cite{yao2017outage,guo2020joint}} &
  \begin{tabular}[c]{@{}c@{}}\cite{khan2019reinforcement,elsayed2020transfer,kwon2019multi}\\ \cite{khosravi2020learning,guo2020joint}\end{tabular} &
  \multicolumn{1}{c|}{\begin{tabular}[c]{@{}c@{}}\cite{khan2019reinforcement,khosravi2020learning,kim2019online,tang2020deep,pateromichelakis2018context,tafintsev2020aerial,shen2020mobility,lei2020deep}\\ \cite{elsayed2020transfer,guo2020joint}\\ \textbf{Ours}\end{tabular}} &
  \multicolumn{1}{c|}{\cite{ortiz2019scaros}} &
  \multicolumn{1}{c|}{\cite{yao2017outage}} &
  \cite{kwon2019multi} \\ \hline
\multicolumn{1}{|c|}{\multirow{5}{*}{SARL}} &
  MAB &
  \multicolumn{1}{c|}{\cite{ortiz2019scaros}} &
  \multicolumn{1}{c|}{\cite{kim2019online,ortiz2019scaros}} &
  \multicolumn{1}{c|}{} &
  \multicolumn{1}{c|}{} &
  \multicolumn{1}{c|}{} &
  \multicolumn{1}{c|}{} &
   &
  \multicolumn{1}{c|}{\cite{kim2019online}} &
  \multicolumn{1}{c|}{\cite{ortiz2019scaros}} &
  \multicolumn{1}{c|}{} &
   \\ \cline{2-13} 
\multicolumn{1}{|c|}{} &
  QL-based &
  \multicolumn{1}{c|}{} &
  \multicolumn{1}{c|}{} &
  \multicolumn{1}{c|}{} &
  \multicolumn{1}{c|}{} &
  \multicolumn{1}{c|}{\cite{khosravi2020learning,elsayed2020transfer}} &
  \multicolumn{1}{c|}{} &
  \cite{elsayed2020transfer,khosravi2020learning} &
  \multicolumn{1}{c|}{\cite{khosravi2020learning,elsayed2020transfer}} &
  \multicolumn{1}{c|}{} &
  \multicolumn{1}{c|}{} &
   \\ \cline{2-13} 
\multicolumn{1}{|c|}{} &
  DQN-based &
  \multicolumn{1}{c|}{} &
  \multicolumn{1}{c|}{} &
  \multicolumn{1}{c|}{\cite{tang2020deep,feng2019dealing}} &
  \multicolumn{1}{c|}{\cite{lei2020deep}} &
  \multicolumn{1}{c|}{} &
  \multicolumn{1}{c|}{} &
   &
  \multicolumn{1}{c|}{\cite{tang2020deep,lei2020deep,feng2019dealing}} &
  \multicolumn{1}{c|}{} &
  \multicolumn{1}{c|}{} &
   \\ \cline{2-13} 
\multicolumn{1}{|c|}{} &
  AC-based &
  \multicolumn{1}{c|}{\cite{zhang2021resource}} &
  \multicolumn{1}{c|}{\cite{zhang2021resource,gupta2020learning}} &
  \multicolumn{1}{c|}{} &
  \multicolumn{1}{c|}{\cite{lei2020deep}} &
  \multicolumn{1}{c|}{} &
  \multicolumn{1}{c|}{} &
   &
  \multicolumn{1}{c|}{\cite{lei2020deep,zhang2021resource}} &
  \multicolumn{1}{c|}{\cite{gupta2020learning}} &
  \multicolumn{1}{c|}{} &
   \\ \cline{2-13} 
\multicolumn{1}{|c|}{} &
  Others &
  \multicolumn{1}{c|}{\cite{vu2018path}} &
  \multicolumn{1}{c|}{} &
  \multicolumn{1}{c|}{\cite{vu2018path}} &
  \multicolumn{1}{c|}{} &
  \multicolumn{1}{c|}{\cite{vu2018ultra}} &
  \multicolumn{1}{c|}{\cite{vu2018ultra}} &
   &
  \multicolumn{1}{c|}{} &
  \multicolumn{1}{c|}{} &
  \multicolumn{1}{c|}{} &
  \cite{vu2018path,vu2018ultra} \\ \hline
\multicolumn{2}{|c|}{MARL} &
  \multicolumn{1}{c|}{\textbf{Ours}} &
  \multicolumn{1}{c|}{\textbf{Ours}} &
  \multicolumn{1}{c|}{} &
  \multicolumn{1}{c|}{\cite{kwon2019multi}} &
  \multicolumn{1}{c|}{} &
  \multicolumn{1}{c|}{\cite{guo2020joint}} &
  \cite{khan2019reinforcement,kwon2019multi,guo2020joint,sana2020multi} &
  \multicolumn{1}{c|}{\begin{tabular}[c]{@{}c@{}}\cite{khan2019reinforcement,guo2020joint,sana2020multi}\\ \textbf{Ours}\end{tabular}} &
  \multicolumn{1}{c|}{} &
  \multicolumn{1}{c|}{} &
  \cite{kwon2019multi} \\ \hline
\end{tabular}%
}
\end{table*}
